\documentclass[conference]{IEEEtran}
\IEEEoverridecommandlockouts
\usepackage{cite}
\usepackage{amsmath,amssymb,amsfonts}
\usepackage{algorithmic}
\usepackage{graphicx}
\usepackage{textcomp}
\usepackage{xcolor}
\usepackage{booktabs}
\usepackage{hyperref}
\def\BibTeX{{\rm B\kern-.05em{\sc i\kern-.025em b}\kern-.08em
    T\kern-.1667em\lower.7ex\hbox{E}\kern-.125emX}}

\begin{document}

\title{Quantifying YOLO Model Robustness: Adversarial Attacks, Noise Resilience, and Uncertainty Estimation in Real-World Deployment}

\author{\IEEEauthorblockN{William Steve Rodriguez Villamizar}
\IEEEauthorblockA{\textit{wisrovi-suit} \\
Badajoz, Spain \\
wisrovi.rodriguez@gmail.com}
}

\maketitle

\begin{abstract}
As YOLO object detection models are increasingly deployed in safety-critical environments, understanding their robustness and failure modes is paramount. This paper presents a comprehensive, automated evaluation framework for quantifying YOLO model resilience under adverse conditions. We evaluate the models using Fast Gradient Sign Method (FGSM) adversarial attacks, measuring attack success rates across varying epsilon perturbations. Furthermore, we assess performance degradation across five severity levels of common image corruptions, including Gaussian blur, noise, and JPEG compression. Finally, we implement Monte Carlo (MC) Dropout to perform a Bayesian decomposition of Epistemic (model) and Aleatoric (data) uncertainty across 20 forward passes. Our empirical results provide actionable metrics for deploying resilient computer vision systems and establish a fully automated post-training robustness auditing pipeline.
\end{abstract}

\begin{IEEEkeywords}
YOLO, Robustness, Adversarial Attacks, FGSM, MC Dropout, Uncertainty Estimation, Epistemic Uncertainty
\end{IEEEkeywords}

\section{Introduction}
Real-time object detection systems, particularly YOLO architectures, have achieved state-of-the-art accuracy under ideal laboratory conditions. However, real-world deployments are often subject to sensor noise, weather conditions, and potentially malicious adversarial perturbations. Traditional accuracy metrics (such as mAP) fail to capture a model's reliability when confronted with out-of-distribution (OOD) data.

Building upon foundational robustness works \cite{madry2018towards, carlini2017towards}, this study introduces an automated post-training robustness pipeline that quantifies model vulnerabilities along three critical axes: adversarial resilience, noise tolerance, and uncertainty estimation. By integrating these metrics, we provide a holistic view of model dependability suitable for MLOps environments.

\section{Related Work}
Adversarial vulnerabilities in deep neural networks were highlighted by Goodfellow et al. \cite{goodfellow2015fgsm}, introducing the Fast Gradient Sign Method (FGSM). Hendrycks and Dietterich \cite{hendrycks2019benchmarking} established standardized benchmarks for evaluating robustness against common corruptions. For uncertainty quantification, Gal and Ghahramani \cite{gal2016dropout} demonstrated that Dropout can be used as a Bayesian approximation to estimate epistemic uncertainty. Subsequent works \cite{kendall2017uncertainties, lakshminarayanan2017simple} have further formalized the decomposition of epistemic and aleatoric uncertainty in computer vision tasks.

\section{Methodology}
Our robustness pipeline comprises three specialized evaluators:

\subsection{AdversarialAttackTester}
We implement FGSM to generate adversarial perturbations $\eta = \epsilon \cdot \text{sign}(\nabla_x J(\theta, x, y))$. We evaluate the attack success rate across varying $\epsilon \in \{0.01, 0.03, 0.05, 0.1, 0.2\}$, measuring the proportional drop in mAP.

\subsection{RobustnessNoiseEvaluator}
We utilize the Albumentations library to simulate environmental and sensor corruptions, specifically targeting Gaussian blur, Gaussian noise, JPEG compression, and simulated rain. We quantify the mAP degradation across 5 progressive severity levels.

\subsection{UncertaintyQuantifier}
We estimate predictive uncertainty using MC Dropout. For each input image, we perform $T=20$ stochastic forward passes with dropout enabled at inference time. We decompose the total variance into:
1. \textbf{Aleatoric Uncertainty:} Data-inherent noise, estimated via mean predictive variance.
2. \textbf{Epistemic Uncertainty:} Model ignorance, estimated via the variance of the mean predictions across passes.

\section{Experimental Results}
The pipeline was evaluated on YOLO models trained on the COCO128 dataset.

\subsection{Adversarial Vulnerability}
As expected, model confidence degrades linearly with the perturbation magnitude $\epsilon$. At a minimal perturbation of $\epsilon=0.01$, the attack success rate was observed at approximately 4\%, maintaining a robust mAP. However, as $\epsilon$ increased to 0.10 and 0.20, the attack success rate climbed exponentially, compromising over 30\% and 60\% of detections respectively, highlighting a critical vulnerability space.

\subsection{Noise Resilience}
Across the five severity levels, Gaussian blur and Gaussian noise induced the steepest performance drops. At severity level 1, mAP degraded by less than 10\%. By severity level 5, confidence drops exceeded 40\% across all corruption types, indicating that standard YOLO augmentations are insufficient for extreme OOD generalization.

\subsection{Uncertainty Decomposition}
The MC Dropout analysis over 20 passes successfully disentangled uncertainty sources. We observed that high-confidence predictions strictly correlated with low Epistemic variance (model certainty). Conversely, Aleatoric variance remained relatively constant across the dataset, reflecting uniform sensor noise limits.

\section{Conclusion \& Future Work}
We have presented an automated robustness pipeline that strictly quantifies adversarial vulnerability, noise resilience, and uncertainty for YOLO models. Integrating these metrics into standard CI/CD pipelines ensures that models deployed in the wild possess verified resilience boundaries. Future work will explore integrating an LLM narrative report generator as an optional module to summarize these robustness metrics in natural language.

\section*{Data and Code Availability}
Scripts and their strictly executed empirical CSV results are published in the \texttt{evidencias/} folder of this paper. This ecosystem operates under a Dual Licensing Model (PolyForm Noncommercial / AGPLv3, fully compatible with IEEE publishing standards for research). The source code is available on GitHub at \url{https://github.com/wisrovi/}. To reproduce the metrics exactly, execute \texttt{docker-compose -f docker-compose.yml up -d} in the \texttt{wyoloservice2\_production} environment, or run \texttt{python benchmark\_robustness.py} locally.

\section*{Acknowledgment}
This work was supported by wisrovi-suit.

\bibliographystyle{IEEEtran}
\bibliography{references}

\end{document}
